\section{Projection and rounding to real signs}\label{sec:completion}

We now complete the proof of Theorem~\ref{thm:main}. The functions in
Proposition~\ref{prop:waves} already have almost constant modulus.
Their Fourier coefficients are real and lie in an almost unit cube
after multiplication by $\sqrt N$. Their large $L^2$ norm forces
most of these coefficients to be close to signs, exactly the condition
needed by Lemma~\ref{lem:complex-rounding}.

Fix a sufficiently small $\delta>0$, and let $B_N$ be supplied by
Proposition~\ref{prop:waves}. Choose an absolute constant $C_1$ so that
\[
 S_\delta=1+C_1\sqrt\delta,
 \qquad Y_k=\frac{\sqrt N}{S_\delta}\widehat B_N(k)
 \quad(0\le k<N)
\]
satisfy $Y_k\in[-1,1]$ for every sufficiently large $N$.
The symmetry $B_N(-t)=\overline{B_N(t)}$ makes every Fourier
coefficient real. Set
\[
 U_Y(t)=\frac1{\sqrt N}\sum_{k=0}^{N-1}Y_k\e(kt).
\]
The absolute Fourier representation and exterior-tail estimate in
Proposition~\ref{prop:waves} give
\begin{equation}\label{eq:projection-uniform}
 U_Y=B_N/S_\delta+O_\delta(N^{-1})
 \quad\text{uniformly on }\T.
\end{equation}
Parseval's identity and $|B_N|\ge1$ consequently imply
\[
 \frac1N\sum_{k=0}^{N-1}Y_k^2\ge S_\delta^{-2}-o(1).
\]
Here and below $N\to\infty$ with all the data chosen from $\delta$
fixed. Since $|Y_k|\ge Y_k^2$, the defect satisfies
\[
 \frac\mu N:=\frac1{2N}\sum_{k=0}^{N-1}(1-|Y_k|)
 \le\tfrac12(1-S_\delta^{-2})+o(1).
\]
Put $q_\delta=\tfrac12(1-S_\delta^{-2})+\delta$.
Then $q_\delta\to0$ as $\delta\downarrow0$, and for all sufficiently
large $N$ we have $\mu/N\le q_\delta<1/2$.

Lemma~\ref{lem:complex-rounding}, applied to these real inputs, produces
$\varepsilon_k\in\{-1,1\}$ such that
\[
 U_\varepsilon(t)=\frac1{\sqrt N}\sum_{k=0}^{N-1}\varepsilon_k\e(kt)
 \quad\text{satisfies}\quad
 \norm{U_\varepsilon-U_Y}_\infty
 \le C\left(N^{-1/2}+
        \sqrt{q_\delta\log(80/q_\delta)}\right).
\]
We used the monotonicity of $q\log(80/q)$ on $(0,1/2]$;
the zero-defect case causes no change. If
$E_\delta=C\sqrt{q_\delta\log(80/q_\delta)}$, the two-sided modulus
bound in Proposition~\ref{prop:waves} and
\eqref{eq:projection-uniform} give, uniformly in $t$,
\[
 S_\delta^{-1}-E_\delta-o(1)
 \le |U_\varepsilon(t)|
 \le \frac{1+C\delta}{S_\delta}+E_\delta+o(1).
\]
For fixed $\delta$, the $o(1)$ terms vanish as $N\to\infty$;
the resulting bounds both tend to $1$ as $\delta\downarrow0$.
Given the $\varepsilon$ of Theorem~\ref{thm:main},
choose $\delta$ so that the limiting bounds lie strictly between
$1-\varepsilon$ and $1+\varepsilon$, and then take $N$ large enough
to absorb the displayed errors. This proves the theorem.

The construction imposes no divisibility condition on $N$.
All torus dimensions, packing parameters, intervals, and leading phases were
fixed before $N$ was chosen. The phase identities use only that the
Fourier index $k$ is an integer. The conclusion therefore holds for
every sufficiently large integer length.
